\bigskip
\bigskip
\subsubsection*{Խնդիրներ և հարցեր թեմա 16-ի վերաբերյալ}

\begin{enumerate}[label=\thesection.\arabic*.]

% 17.1
\item Դիցուք $(X, \tau_1)$ և $(X, \tau_2)$ տոպոլոգիական տարածություններն այնպիսին են, որ $\tau_1 \subset \tau_2$։ Ճի՞շտ է պնդումը.
\begin{enumerate}
    \item[ա)] եթե $(X, \tau_1)$-ը կոմպակտ է, ապա $(X, \tau_2)$-ը կոմպակտ է;
    \item[բ)] եթե $(X, \tau_2)$-ը կոմպակտ է, ապա $(X, \tau_1)$-ը կոմպակտ է:
\end{enumerate}

% 17.2
\item Կոմպա՞կտ է արդյոք աջից կիսաբաց ինտերվալների $(\R, \mapsto)$ տարածությունը:

% 17.3
\item Կոմպա՞կտ է արդյոք ա) $[1, 2)$; բ) $(1, 2]$ ենթաբազմությունը $(\R, \mapsto)$ տարա\-ծու\-թյու\-նում:

% 17.4
\item Գոյություն ունի՞ $(\R, \text{սովոր.})$ տարածության
\begin{enumerate}
    \item[ա)] $(-1, 1)$ ինտերվալի;
    \item[բ)] $[-1, 1]$ հատվածի
\end{enumerate}
անընդհատ արտապատկերում $\R$-ի վրա:

% 17.5
\item Ապացուցեք․
\begin{enumerate}
    \item[ա)] էվկլիդեսյան $\R^{n+1}$ տարածության $S^n$ միավոր սֆերան և $B^{n+1}$ միավոր փակ գունդը կոմպակտ տարածություններ են;
    \item[բ)] $\mathbb{RP}^n$, $n \ge 0$ պրոյեկտիվ տարածությունները կոմպակտ տարածու\-թյուն\-ներ են:
\end{enumerate}

% 17.6
\item Ապացուցեք․ տոպոլոգիական տարածության ամեն երկու կոմպակտ ենթա\-բազ\-մու\-թյուն\-նե\-րի միավորումը կոմպակտ ենթաբազմություն է:

% 17.7
\item Ի տարբերություն նախորդ խնդրի պնդման՝ երկու կոմպակտ ենթաբազմու\-թ\-յուն\-ների հատումը կարող է չլինել կոմպակտ ենթաբազմություն:
\par Բերենք օրինակ․ դիտարկենք իրական թվերի $A$ ենթաբազմությունը՝ կազմված բոլոր ամբողջ թվերից և $\pm \dfrac{1}{2}$ թվերից՝ $A = \Z \cup \left\{\pm \dfrac{1}{2}\right\}$։ Սահմանենք տոպոլո\-գիա $A$ բազմության վրա՝ համարելով բաց ենթաբազմություններ $\varnothing$-ը և $A$-ն, $\Z$-ի բոլոր ենթաբազմությունները և $A$-ի այն բոլոր ենթաբազմությունները, որոնց լրացումները վերջավոր են։ Ցույց տվեք․
\begin{enumerate}
    \item[1.] տոպոլոգիայի 1-3 աքսիոմները բավարարվում են;
    \item[2.] $A$-ի $B = \Z \cup \left\{-\dfrac{1}{2}\right\}$ և $C = \Z \cup \left\{\dfrac{1}{2}\right\}$ ենթաբազմությունները կոմպակտ են;
    \item[3.] $B \cap C$ ենթաբազմությունը կոմպակտ չէ:
\end{enumerate}

% 17.8
\item Ճի՞շտ է պնդումը․ ցանկացած մետրիկային տարածության ամեն մի փակ և սահմանափակ ենթաբազմություն կոմպակտ ենթաբազմություն է:

\begin{hint}
Որպես ժխտօրինակ դիտարկեք որևէ $X$ անվերջ բազմություն 
\[
\rho(x, y) = \begin{cases}
    0, & \text{երբ } x=y \\
    1, & \text{երբ } x\ne y 
\end{cases}
\]
մետրիկով և նրա $\mathcal{D}(x_0, 1) = \{x;\ \rho(x_0, x) \le 1 \}$ ենթաբազմությունը, որտեղ $x_0$-ն $X$-ի որևէ սևեռված կետ է:
\end{hint}

% 17.9
\item Ապացուցեք․ եթե $A$-ն $X$ հաուսդորֆյան տարածության կոմպակտ ենթա\-բազ\-մություն է, ապա $A$-ի և ցանկացած $x_0 \in X \setminus A$ կետի համար գոյություն ունեն չհատվող շրջակայքեր:

\begin{hint}
Օգտվեք \hyperref[ch16thm5]{թեորեմ 5}-ի ապացուցման ընթացքից:
\end{hint}

% 17.10
\item Ապացուցեք․ ամեն մի հաուսդորֆյան կոմպակտ տարածություն ռեգուլյար տարա\-ծություն է:

% 17.11
\item Ապացուցեք․ հաուսդորֆյան տարածության ցանկացած երկու կոմպակտ չհատ\-վող ենթաբազմությունների համար գոյություն ունեն դրանց չհատվող շրջա\-կայ\-քեր:

% 17.12
\item Ապացուցեք․ ամեն մի հաուսդորֆյան կոմպակտ տարածություն նորմալ տարա\-ծություն է:

\end{enumerate}
